\documentclass[a4paper,slidestop,xcolor=pst,dvips,blue]{beamer}

\usepackage{beamerthemesplit}
\usepackage[utf8]{inputenc}
\usepackage[spanish]{babel}
\usepackage{graphicx}
\usepackage{pstricks} % PSTricks package
\usepackage{setspace}
\usepackage{multirow}
\usepackage{listings}
\usepackage{pgfpages}
\usepackage{hyperref}
\usepackage{etoolbox}
\usepackage{epstopdf}

\makeatletter
\patchcmd{\beamer@sectionintoc}{\vskip1.5em}{\vskip0.5em}{}{}
\makeatother

\setbeamercovered{dynamic}
\setcounter{tocdepth}{2}
\setbeamercolor{frametitle}{fg=black,bg=white}
\setbeamercolor{section in toc shaded}{fg=black}
\setbeamercolor{section in toc}{fg=red}
\setbeamercolor{subsection in toc shaded}{fg=black}
\setbeamercolor{subsection in toc}{fg=red}
\setbeamerfont{section in toc}{size=\small}
\setbeamerfont{subsection in toc}{size=\small}
\setbeamertemplate{section in toc shaded}[default][99]
\setbeamertemplate{subsection in toc shaded}[default][99]

\AtBeginSection[]
{\begin{frame}[c]
  \frametitle{Índice}
	\tableofcontents[currentsection,
        sectionstyle=show/shaded,
        subsectionstyle=hide]
\end{frame}}

\AtBeginSubsection[]
{\begin{frame}[c]
	\frametitle{Índice}
	\tableofcontents[
  		currentsection,
  		sectionstyle=shaded/shaded,
  		currentsubsection,
  		subsectionstyle=show/shaded/hide
		]
\end{frame}}

\setbeamercolor{frametitle}{fg=black,bg=white}

\setbeamertemplate{frametitle}{
	\begin{centering}
		\insertframetitle
		\par
	\end{centering}
}

\usetheme[secheader]{Boadilla} 

\usetheme[secheader]{Boadilla}

\title[Procesos Ing. Requisitos]{Procesos de Ingeniería de Requisitos}

\author[Pablo Sánchez]{\alert{Pablo Sánchez}}

\institute[I2E]{
		   Dpto. Ingeniería Informática y Electrónica \\
		   Universidad de Cantabria \\
		   Santander (Cantabria, España) \\
		   \texttt{p.sanchez@unican.es}
}

\date{}

\begin{document}

\begin{frame}[c]
	\titlepage
	\begin{columns}
		\column{0.50\linewidth}
			\centering
			\includegraphics[width=.28\textwidth,keepaspectratio=true]{images/istr.eps}
		\column{0.50\linewidth}
			\centering
			\includegraphics[width=.25\textwidth,keepaspectratio=true]{images/uc.eps}
    \end{columns}
\end{frame}

\begin{frame}[c]
    \frametitle{\alert{Advertencia}}
    \begin{center}
        Todo el material contenido en este documento  no constituye en modo alguno una obra de referencia o apuntes oficiales mediante el cual se puedan preparar las pruebas evaluables necesarias para superar la asignatura de Ingeniería de Requisitos. \ \\
        \ \\
        Este documento contiene exclusivamente una serie de diapositivas cuyo objetivo es servir de complemento visual a las actividades realizadas en el aula para la transmisión del contenido sobre el cual versarán las mencionadas pruebas evaluables.  \ \\
        \ \\
        Dicho de forma más clara, \alert{estas transparencias no son apuntes y su objetivo no es en modo alguno servir para que el alumno pueda preparar la asignatura.}
    \end{center}
\end{frame}

\section{Introducción}

\begin{frame}[c]
    \frametitle{Objetivos del Tema}
    \begin{enumerate}[<+->]
         \item Conocer y comprender la importancia y objetivo último de la Ingeniería de Requisitos.
         \item Conocer y comprender la terminología básica general relacionada con la Ingeniería de Requisitos.
         \item Conocer y comprender cómo funcionan los procesos de Ingeniería de Requisitos.
         \item Conocer y comprender cómo se integra la Ingeniería de Requisitos en los procesos de desarrollo software.
    \end{enumerate}
\end{frame}

\begin{frame}[c]
    \frametitle{Bibliografía}
    \begin{thebibliography}{1}
        \bibitem{pohl:interno} Klaus Pohl. \newblock {\em {Requirements Engineering: Fundamentals, Principles and
        Techniques}}.\newblock Springer, June 2010.
    \end{thebibliography}
\end{frame}

\section{Objetivos y Necesidad de la Ingeniería de Requisitos}

\begin{frame}
    \frametitle{¿Por Qué la Ingeniería de Requisitos ?}
    \begin{center}
        \includegraphics[width=\linewidth]{images/introduction/ciclo00.eps}
    \end{center}
\end{frame}

\begin{frame}[c]
    \frametitle{¿Por Qué Ingeniería de Requisitos ?}
    \begin{enumerate}[<+->]
        \item Una inapropiada Ingeniería de Requisitos es la causa principal del fracaso de un proyecto sw.
            %% NOTA: \emph{Standish Group Report 2011}: 63\% Proyectos Fallidos, 48\% Fallos debido a mala Ing. Requisitos
            %% NOTA: Sistema de ambulancias de Londres.
        \item Una apropiada Ingeniería de Requisitos reduce el número de defectos y el coste asociado a su reparación.
            %% Costes error no detectado en Ing. Requisitos x20
    \end{enumerate}
\end{frame}

\section{Terminología Básica}

\subsection{Requisito Sw}

\begin{frame}
    \frametitle{Definición de Requisito}
    \begin{block}{Requisito Sw~\cite{ieee:seGlossary}}
        Un \alert{\emph{requisito software}} se define como:
        \begin{enumerate}
            \item<2-> Capacidad o condición que es precisada por un usuario para resolver un problema o satisfacer un objetivo.
            \item<3-> Capacidad o condición que debe satisfacer o poseer un sistema, un componente del mismo, de acuerdo a un contrato, estándar, especificación u cualquier otro documento formalmente definido.
            \item<4-> Representación documentada de una condición o capacidad, de acuerdo a las definiciones (1) o (2).
        \end{enumerate}
    \end{block}
    \uncover<5->{
        \begin{block}{Requisito Sw~\cite{ieee:swebok}}
            Un \alert{\emph{requisito software}} es una propiedad que debe poseer un sistema software desarrollado o adaptado para solventar un determinado problema.
        \end{block}
    }
\end{frame}

\subsection{Requisito Producto vs Requisito Proceso}

\begin{frame}
    \frametitle{Producto vs Proceso}
    \begin{block}{Requisito de Producto~\cite{ieee:swebok}}
            Una \alert{\emph{requisito (sw) de producto}} es un requisito sw que se refiere a las propiedades que el producto sw a ser desarrollado debe exhibir.
    \end{block}
    \uncover<2->{
        \begin{block}{Requisito de Proceso~\cite{ieee:swebok}}
            Una \alert{\emph{requisito (sw) de proceso}} es un requisito que se refiere a las propiedades que debe exhibir el proceso utilizado para el desarrollo de un cierto sistema sw.
        \end{block}
    }
\end{frame}

\subsection{Requisitos Funcionales y No Funcionales}

\begin{frame}
    \frametitle{Requisitos Funcionales, No Funcionales y Restricciones}
    \begin{block}{Requisito Funcional~\cite{ieee:swebok}}
        Un \alert{\emph{requisito (sw) funcional}} es un requisito que describe las funciones que el software debe ejecutar.
    \end{block}
    \uncover<2->{
        \begin{block}{Restricción~\cite{robertson:2012}}
            Una \alert{\emph{restricción}} es un requisito a nivel corporativo o tecnológico que restringe el modo en el cual un sistema sw puede ser construido.
        \end{block}
    }
    \uncover<3->{
        \begin{block}{Atributo de Calidad~\cite{pohl:2010}}
            Un \alert{\emph{atributo de calidad (sw)}}, o requisito de calidad, de un sistema sw define una propiedad relativa a la calidad final de un producto sw como un todo, o a la calidad final de alguno de sus componentes, servicios o funciones.
        \end{block}
    }
\end{frame}

\begin{frame}
    \frametitle{Requisitos Funcionales, No Funcionales y Restricciones}
    \begin{block}{Requisito No Funcional~\cite{chung:1999}}
        Un \alert{Requisito No Funcional} de un sistema sw es un requisito que no indica qué debe hacer el sistema, sino cómo debe hacerlo.
    \end{block}
\end{frame}

\subsection{Esencia y Encarnación}

\begin{frame}
    \frametitle{Esencia vs Encarnación}
    \begin{block}{Esencia de un Sistema~\cite{mcmenamin:1984}}
        La \alert{\emph{esencia de un sistema sw}} es el conjunto de requisitos sw que indican las capacidades o caracterásticas esenciales  que debe poseer un sistema sw para satisfacer su objetivo, con independencia de como se implemente dicho sistema.
    \end{block}
    \uncover<2-> {
        \begin{block}{Encarnación de un Sistema~\cite{mcmenamin:1984}}
            La \alert{\emph{encarnación de un sistema sw}} es todo aquello que es necesario para la implementación y funcionamiento de un sistema sw, desde su concepción, hasta el fin de su vida útil.
        \end{block}
    }
\end{frame}

\subsection{Objetivos y Escenarios}

\begin{frame}
    \frametitle{Objetivos y Escenarios}
    \begin{block}{Objetivos~\cite{lamsweerde:2001}}
        Un \alert{\emph{objetivo}} de un sistema sw es un fin que dicho sistema software debe alcanzar.
    \end{block}
    \uncover<2-> {
        \begin{block}{Escenario~\cite{pohl:2010}}
            Un \alert{\emph{escenario}} describe un ejemplo concreto de satisfacción o fallo en la satisfacción de un determinado objetivo o conjunto de objetivos.
        \end{block}
    }
\end{frame}

\section{Procesos de Ingeniería de Requisitos}

\begin{frame}[c]
    \frametitle{Procesos de Ingeniería de Requisitos}
    \begin{center}
		\includegraphics[width=\textwidth,keepaspectratio=true]{images/framework/ref.eps}
    \end{center}
\end{frame}

\section{IR en los Procesos de Desarrollo}

\begin{frame}[c]
    \frametitle{Ing. de Requisitos y Procesos de Negocio}
    \begin{center}
			\includegraphics[width=\textwidth,keepaspectratio=true]{images/introduction/ingReqOrganizacion.eps}
    \end{center}
\end{frame}

\begin{frame}
    \frametitle{Ing. de Requisitos y Procesos de Desarrollo}
    \begin{center}
			\includegraphics[width=0.80\textwidth,keepaspectratio=true]{images/introduction/ingReqDesarrollo.eps}
    \end{center}
\end{frame}

\begin{frame}
    \frametitle{Modelo \emph{Twin Peaks}}
    \begin{center}
		\includegraphics[width=.9\textwidth,keepaspectratio=true]{images/introduction/twinPeaksModel.eps}
    \end{center}
\end{frame}

\section{Sumario y Referencias}

\subsection{Sumario}

\begin{frame}
    \frametitle{Objetivo Ingeniería de Requisitos}
    \begin{block}{Ingeniería de Requisitos}
        La \alert{\emph{Ingeniería de Requisitos}} es un proceso cooperativo, iterativo e incremental cuyo objetivo es el de asegurar que:
        \begin{enumerate}
            \item<2-> Se conocen y comprenden, con el nivel de detalle suficiente, todos los requisitos que son de relevancia para el desarrollo del sistema.
            \item<3-> Existe un consenso entre los diferentes actores con intereses en el sistema acerca de dichos requisitos.
            \item<4-> Todos estos requisitos están documentados y descritos conforme a formatos y normativas establecidas.
        \end{enumerate}
    \end{block}
    \uncover<5->{
        \begin{block}{Objetivo de la Ingeniería de Requisitos}
            El objetivo de la Ingeniería de Requisitos es \emph{establecer una visión en un contexto}~\cite{jarke:1993}.
        \end{block}
    }
\end{frame}

\begin{frame}
    \frametitle{¿Qué Tengo que Saber de Todo Esto?}
    \begin{enumerate}[<+->]
        \item Conocer y comprender la utilidad de la Ingeniería de Requisitos.
        \item Conocer y comprender la terminología básica asociada a la Ingeniería de Requisitos.
        \item Conocer y comprender el funcionamiento y etapas de los procesos de Ingeniería de Requisitos
        \item Conocer y comprender el papel de la Ingeniería de Requisitos en los procesos de desarrollo sw.
    \end{enumerate}
\end{frame}

\subsection{Referencias}

\begin{frame}[allowframebreaks]
	\frametitle{Referencias}
	\bibliographystyle{apalike}
	\bibliography{ir}
\end{frame}

\end{document}
